\section{Temporal matching of product slice corners}\label{sec:temporal}
We now restrict to $D=3$, $4\mid n$, $n\geq32$, and the orthogonal
internal basis. Thus $a=c/48$ and $H=cn^2$. This section concerns the actual
coexact domain $U=U_0^2$ and the original mean normalization, but permits
an enlarged observation class: cyclic intervals may wrap in the fixed
chart. We call this the cyclic extension. Section 1's canonical class
uses nonwrapping chart intervals and is not redefined here. Every
nonwrapping product-slice partition is a member of this extension, so
restricting the theorem to those partitions gives a canonical necessary
bound. The rapid example below belongs only to the cyclic extension.
This supplies no unrestricted lower bound for every canonical partition.

At each longitudinal slice $x_1=i$, let $\{J_{i,j}\}_j$ and
$\{K_{i,k}\}_k$ be nonempty cyclic interval partitions of $\Z_n$ whose
product of factor counts is the same integer $q$ for every $i$. The cells
of channel 1 are $\{i\}\times J_{i,j}\times K_{i,k}$. Their full spatial-row
count is $3nq$. The second channel can use any permitted partition. Write
$E_i$ for the product of the terminal endpoints of the two interval
partitions, so that $E_i$ has exactly $q$ distinct transverse points.
For $K\geq1$, let $S_K$ be the maximum over cyclic $i$ of the minimum
bottleneck matching distance between $E_i$ and $E_{i+K}$, using the
periodic transverse $\ell^1$ metric. These extrema exist over finite sets.
We impose no common labelling around the cycle.

\begin{theorem}\label{thm:temporal}
Put $L=\log((n/4+1)/3)$. If for some integer $K\geq1$,
\begin{equation}\label{eq:windowthreshold}
 L\geq 2^{18}q A_K,\qquad
 A_K=\max\{16K,9K+6S_K+S_K^2\},
\end{equation}
there exists a real nonzero $z\in U$ with every original mean zero and
$Q(z)<7c\norm z^2/16$. In particular, no finite quadratic penalty on those
means can restore target $c$. Consequently recovery requires
\begin{equation}\label{eq:windownecessary}
 L<2^{18}q\max\{16K,9K+6S_K+S_K^2\}
 \quad\hbox{for every integer }K\geq1.
\end{equation}
The assertion permits $K\geq n$ and includes the longitudinal seam.
\end{theorem}

We give the constructive proof, including the norm and original-source
energy bounds. No assumed capacity estimate, projection invariance or
coexact gap is used as an input.

\subsection{A finite transverse cutoff}
Let $Y=\Z_n^2$, $R=n/4$, and let $\rho_p$ be periodic $\ell^\infty$
distance from $p\in Y$. Set
\[
 f_p(y)=\min\{1,\max\{0,L^{-1}\log((\rho_p(y)+1)/3)\}\}.
\]
It is zero for $\rho_p\leq2$ and one for $\rho_p\geq R$. We claim
\begin{equation}\label{eq:radial}
 \sum_{y,h=2,3}|\partial_h f_p(y)|^2\leq\frac{64}{L},
 \qquad \norm{1-f_p}_Y^2\leq\frac{4n^2}{L^2}.
\end{equation}
For a nonconstant edge, the smaller radius is $2\leq r<R$. There are
$8r+4\leq32r$ forward-labelled edges crossing that square boundary,
disregarding direction. Since $r+1\leq R\leq n/4$, no antipodal ambiguity
enters this count. Each scalar increment is at most $1/((r+1)L)$.
Thus the energy is at most
\[
 \frac{32}{L^2}\sum_{r=2}^{R-1}\frac{r}{(r+1)^2}
 \leq\frac{32\log(R+1)}{L^2}\leq\frac{64}{L},
\]
where $R\geq8$ gives $\log(R+1)\leq2L$.
For the norm, group sites by
$(R+1)/2^{j+1}<\rho_p+1\leq(R+1)/2^j$, $j\geq0$.
Each group has at most $4(R+1)^2/4^j$ sites and deficit at most
$(j+1)\log2/L$. The latter estimate holds also where clipping makes
$f_p=0$. Since $\sum_{j\geq0}(j+1)^2/4^j=80/27<4$ and $R+1\leq n/2$,
the deficit norm is at most $16(R+1)^2/L^2\leq4n^2/L^2$.
This proves~\eqref{eq:radial}. Telescoping a shortest lattice path also gives
\begin{equation}\label{eq:singletranslation}
 \norm{f_p(\cdot+u)-f_{p'}}_Y
 \leq (|u|_1+d(p,p'))\sqrt{64/L}.
\end{equation}
Here $|u|_1$ can be any chosen integer representative length, and all seams
are periodic.

\subsection{Tagged windows and their translation cost}
Fix $K$ and write $S=S_K$. At slice $i$ take the tagged list
$F_i=\{(j,p):0\leq j<K,\ p\in E_{i+j}\}$. Repeated spatial points keep
their tags, so its size is exactly $Kq$ even if $K\geq n$.
Match $F_i$ to $F_{i+1}$ by keeping the designated $(K-1)q$ shared entries
at the same spatial position and matching the remaining $q$ outgoing
entries to $E_{i+K}$ with distance at most $S$. Some replaced entries can
also be stationary; no claim about the total number of stationary matches
is needed. Define
\[
 \eta(i,y)=\min_{(j,p)\in F_i} f_p(y).
\]
It lies in $[0,1]$ and vanishes within radius 2 of $E_i$. Squared maximum
is bounded by the sum of squares, so~\eqref{eq:radial} gives
\begin{equation}\label{eq:etadeficit}
 \norm{1-\eta}_X^2\leq 4KqV/L^2.
\end{equation}
For two equally indexed real lists, $|\min a-\min b|\leq\max|a-b|$.
Using the designated matching and~\eqref{eq:singletranslation},
\begin{align}\label{eq:movingtranslation}
 \norm{\eta(i+1,\cdot+u)-\eta(i,\cdot)}_Y^2
 &\leq\frac{64q}{L}\{(K-1)|u|_1^2+(|u|_1+S)^2\}\notag\\
 &=\frac{64q}{L}(K|u|_1^2+2|u|_1S+S^2).
\end{align}
With no longitudinal shift, identity matching of all $Kq$ entries instead
gives $64Kq|u|_1^2/L$. Summing either bound over $i$ adds a factor $n$.
The window is only a cutoff recipe; the observed cells remain the original
$q$ cells per slice.

\subsection{Actual coexact membership and observation zeros}
Put $\delta_h f(x)=f(x)-f(x-e_h)$. For a real periodic $\Psi$, take
$\phi=\delta_2\delta_3\Psi$ and the actual two-form potential
$\psi_{23}=\phi$, $\psi_{12}=\psi_{13}=0$. Since
$\partial_h^*=-\delta_h$, direct adjunction yields
\[
 d_1^*\psi=(0,\delta_2\delta_3^2\Psi,-\delta_2^2\delta_3\Psi)=P\Psi.
\]
Every component has zero full sum by telescoping. Summing the second
component on a cell telescopes to four corner values of $\delta_3\Psi$;
the third gives the negative four corner values of $\delta_2\Psi$.
Predecessors of initial endpoints are terminal endpoints of the preceding
interval. Hence all relevant corners are in $E_i$, including wrap cells.
A whole-cycle factor cancels identically.

Use the period-four wave
$\Psi_0(x)=\sin(\pi(x_1+x_2+x_3)/2)$, and set
$\Psi=\eta\Psi_0$. At every corner, $\eta$ and its backward transverse
neighbors vanish. Thus both required differences of $\Psi$ vanish there.
Every channel 1 mean of $P\Psi$ is exactly zero, with its original positive
normalization. Take channel 2 identically zero. This constructs an actual
real $z=(P\Psi,0)\in U$ annihilating all observed rows.

\subsection{Nonzero norm and original energy}
Direct period-four differences give
\[z_0=P\Psi_0=(0,w,-w),\qquad w=(2,-2,-2,2).\]
Thus $|z_0(x)|^2=8$,
$\norm{z_0}^2=8V$, and at every original odd source all $\partial_hz_0$
vanish, including across seams because $4\mid n$.
For a finite convolution $\mathcal T=\sum_r a_r\tau^r$, where $\tau^r$
is translation by $r$, subtraction gives
\[
 \mathcal T(\eta\Psi_0)-\eta\mathcal T\Psi_0
 =\sum_r a_r(\tau^r\eta-\eta)\tau^r\Psi_0.
\]
Each nonzero component of $P$ has coefficient absolute sum at most 8,
longitudinal shift 0 and transverse radius at most 3. Since $|\Psi_0|\leq1$,
the no-longitudinal translation bound yields
\begin{equation}\label{eq:pnormerror}
 \norm{P(\eta\Psi_0)-\eta z_0}^2\leq73728nKq/L.
\end{equation}
Each component of $\partial_2P$ or $\partial_3P$ has coefficient absolute
sum at most 16 and transverse radius at most 4. Its squared translation
factor is at most $16K$. Each component of $\partial_1P$ has coefficient
absolute sum at most 16, longitudinal shifts 0 or 1 and transverse radius
at most 3. Its factors are bounded by $9K$ and $9K+6S+S^2$ respectively.
Consequently each derivative commutator has norm at most
$16\sqrt{64nqA_K/L}$. Restricting to original odd sources where
$\eta\partial_hz_0=0$, and summing the six squared component bounds,
\begin{equation}\label{eq:strongerror}
 \sum_{x\text{ odd},h}|\partial_h(P\Psi)(x)|^2
 \leq98304nqA_K/L.
\end{equation}
These are the original strong sources; no weight has been shifted.

By~\eqref{eq:etadeficit}, \eqref{eq:pnormerror} and $n\leq V$,
\[
 \norm{P\Psi-z_0}^2\leq
 \left(64Kq/L^2+147456Kq/L\right)V.
\]
The threshold and $A_K\geq16K$ imply $L\geq2^{22}Kq$. The coefficient
above is at most $2^{-38}+9/256<2$. Reverse triangle therefore gives
$\norm{P\Psi}>\sqrt{8V}-\sqrt{2V}=\sqrt{2V}$.
In particular the constructed field, not just its potential, is nonzero.
Multiplying~\eqref{eq:strongerror} by $H=cn^2$ gives
\[
 Q_{\rm strong}(z)\leq98304cVqA_K/L
 \leq3cV/8<3c\norm z^2/16.
\]
Every directional forward difference has norm at most $2\norm z$, so
$Q_{\rm weak}(z)\leq12a\norm z^2=c\norm z^2/4$.
Adding proves Theorem~\ref{thm:temporal}. Any quadratic penalty vanishes
on the identically zero observation vector.

\subsection{Count consequences and a rapidly moving grid}
For fixed cuts, $S_1=0$ and $A_1=16$. With bounded one-step corner
motion, $A_1$ stays bounded as well. In either class,
\eqref{eq:windownecessary} forces $q$ to grow at least proportionally to
$L$, and hence requires order $n\log n$ full channel 1 rows. This conclusion
is restricted to these product classes, including their nonwrapping
canonical subclasses. Arbitrary canonical rectangle partitions and
rapidly moving product partitions remain outside this count conclusion.

In the cyclic extension, large adjacent motion does not itself guarantee
recovery. For
$n=2^m$, $m\geq5$, let $\kappa=n/4+2$ and $t_i=\kappa i\bmod n$.
Use the singleton $\{t_i\}$ and its cyclic complement for the first
transverse factor, and $\{0\}$ and its complement for the second.
Use these partitions in both original channels. There are $24n$ full rows
and $E_i=\{t_i,t_i-1\}\times\{0,n-1\}$, so $q=4$.
The identity $4\kappa=n+8$ gives $S_4\leq8$ and $A_4\leq148$.
For $L\geq2^{20}\cdot148$, the theorem gives an invisible real nonzero
coexact field below target. Such dyadic sizes exist, but the threshold is
enormous and is not a practical crossover estimate.

The adjacent bottleneck is exactly $\kappa$: the next corner with first
transverse coordinate $t_i+\kappa$ is at least that far from each old
corner, since $\kappa<n/2-1$; translation realizes that maximum distance.
Also $\gcd(\kappa,n)=2$, so the orbit visits $n/2$ positions.
This is a rapidly moving cyclic family with a short resonance. For an
interior singleton, its cyclic complement wraps and is not one of the
nonwrapping chart intervals in Section 1. Splitting a wrapped cell into
chart rectangles introduces separate means; invisibility of their sum
does not imply invisibility of each refined mean. Thus neither this
$24n$ count nor its failure is transferred to that chart refinement.
The theorem does not assert that every rapidly moving $O(n)$ family fails.


\subsection{A separate counterexample for nonwrapping chart cells}
The previous cyclic example does not imply failure of its chart refinement.
We now construct a new cutoff that annihilates every refined mean separately.
This gives a counterexample inside the original nonwrapping canonical class.

\begin{proposition}\label{prop:chartrefinement}
Let $n=2^m$, $m\geq5$, $\kappa=n/4+2$ and $t_i=\kappa i\bmod n$.
At slice $i$ partition the second-coordinate chart into the nonempty members
of $[0,t_i-1]$, $\{t_i\}$ and $[t_i+1,n-1]$, and the third coordinate
into $\{0\}$ and $[1,n-1]$. Use these nonwrapping product cells in both
original channels, with first coordinate $\{i\}$. Their exact full count is
$36n-24$. If
\begin{equation}\label{eq:chartthreshold}
 L\geq2^{18}\cdot6\cdot148,
\end{equation}
there is a real nonzero $z\in U$ with each of these original means zero and
$Q(z)<7c\norm z^2/16$. This particular canonical family therefore cannot
restore target $c$ with any finite quadratic observation penalty.
\end{proposition}
\begin{proof}
Since $\gcd(\kappa,n)=2$, the orbit visits each even position twice.
Thus $t_i=0$ occurs twice, $t_i=n-1$ never occurs, and all other $t_i$ are
interior. There are four cells per channel at the two zero slices and six
at each other slice. Counting all spatial rows in both channels gives
$2\cdot3\{2\cdot4+(n-2)\cdot6\}=36n-24$.

For interior $t_i$, the actual terminal second-coordinate endpoints are
$t_i-1,t_i,n-1$; for $t_i=0$ they are $0,n-1$. The third-coordinate
endpoints are $0,n-1$. Let $G_i$ be the six-entry tagged list obtained
from $t_i,t_i-1\bmod n,n-1$, each paired with those two third-coordinate
endpoints. It covers every actual corner, including the extra chart
boundary. At $t_i=0$ some positions coincide but retain their tags.
Six is a covering-list count, not an assertion of six cells at every slice.

The identity $4\kappa=n+8$ gives a same-tag matching from $G_i$ to
$G_{i+4}$: four entries translate by $(8,0)$ and two are stationary.
Form the 24-entry window $F_i$ from $G_i,G_{i+1},G_{i+2},G_{i+3}$ and
let $\eta(i,y)$ be the minimum of the associated radial functions $f_p$.
The designated 18 shared entries remain stationary between consecutive
windows; the six replaced entries have displacement at most8. Directly
using~\eqref{eq:radial} and the minimum inequality gives
\begin{align*}
 \norm{1-\eta}_X^2&\leq4\cdot4\cdot6V/L^2,\\
 \norm{\eta(i+1,\cdot+u)-\eta(i,\cdot)}_Y^2
 &\leq\frac{64\cdot6}{L}
       \{4|u|_1^2+16|u|_1+64\}.
\end{align*}
Without longitudinal shift the bound is $64\cdot4\cdot6|u|_1^2/L$.
These explicit estimates remain valid with repeated centres. We have not
applied the constant-cell-count theorem to this varying four/six-cell map.

Set $\Psi=\eta\Psi_0$ and use the actual potential and field
$z=(P\Psi,0)$ from the preceding proof. At every actual refined corner,
$\eta$ and its backward transverse neighbours vanish. For each individual
chart interval, its initial-endpoint predecessor is a terminal endpoint of
the preceding interval, with $-1=n-1$ at the chart boundary. Every one of
the four corners in the telescoped cell sum is therefore covered by $G_i$.
Both needed first differences of $\Psi$ vanish there, so each separately
normalized refined mean is zero. No zero sum on a union is used to infer
zero sums on its pieces. The actual curl-adjoint formula proves real
membership in $U$.

For the convolution shifts used by $P$ and its derivatives, the translation
factors are bounded by $A=\max\{64,36+48+64\}=148$. The same coefficient
sums, now applied to the explicit six-centre estimates above, give
\begin{align*}
 \norm{P(\eta\Psi_0)-\eta z_0}^2
    &\leq73728n\cdot4\cdot6/L,\\
 \sum_{x\text{ odd},h}|\partial_h(P\Psi)(x)|^2
    &\leq98304n\cdot6\cdot148/L,\\
 \norm{P\Psi-z_0}^2
    &\leq(64\cdot24/L^2+147456\cdot24/L)V.
\end{align*}
Threshold~\eqref{eq:chartthreshold} implies $L\geq2^{22}\cdot24$.
The last coefficient is at most $2^{-38}+9/256<2$, so
$\norm{P\Psi}>\sqrt{2V}$. The field itself is nonzero. Multiplying the
strong estimate by the original $H=cn^2$ bounds strong energy by
$3cV/8<3c\norm z^2/16$; weak energy is at most
$12a\norm z^2=c\norm z^2/4$. This proves the strict $7c/16$ bound.
All observation penalties vanish on the zero mean vector.
\end{proof}

This is a newly constructed invisible field for this particular
nonwrapping refinement. It does not retroactively validate a split-cell
inference for the old cyclic field. The threshold is enormous and
unoptimized, and no general impossibility of canonical $O(n)$ partitions
or global canonical logarithmic lower bound follows.


